\begin{claim}{G.7}{Proved}{The full sphere localizer is accessible in ZFC}
The null classes of G.5 and Part A satisfy
\[
 \widehat{\mathcal N}_2
       =\widehat{\mathcal N}_2^{\aleph_1},
 \qquad \Nul=\Nul^{\aleph_1}.
\]
Consequently both $\widehat W_2$ and $W_S$ are accessible reversible fundamental localizers in ZFC. A generating set for each is given by the category-of-simplices models of the maps between countable CW complexes whose cofibers lie in its displayed null class. In particular $W_S$ itself is the accessible counterexample asked for in Part B, and the pointed localization derivator in Part G exists in the stipulated small universe.
\end{claim}
\begin{proof}
We give a strict compact model and a compactness argument for small subcomplexes. The proof does not identify the homotopy category of an arrow category with the arrow category of a homotopy category.

\emph{Compact Kan models.}
Let $K=\Delta^n/\partial\Delta^n$, $n\geq2$, and let $R$ be the reduced free simplicial $\FF_2$-module functor, regarded as a monad on pointed simplicial sets. The standard Bousfield--Kan resolution has cosimplicial terms
\[
 X^q=R^{q+1}K\qquad(q\geq0).
\]
The precise construction and its identification with $\FF_2$-completion are in \cite[\S7.7, pp.~1032--1033]{FGHBousfield2003}. Each simplicial degree of each $X^q$ is finite, since $K$ has finitely many simplices in each degree and a free vector space on a finite set over $\FF_2$ is finite. Also each $X^q$ is a Kan complex, being a simplicial abelian group.

Here is a direct check that this cosimplicial resolution is Reedy fibrant. Its codegeneracies $s^i$ are linear for the outer $R$-module structures, as are all cofaces $d^i$ with $i\geq1$. A matching tuple in degree $q$ is determined by compatible elements $(a_0,\ldots,a_{q-1})$ in degree $q-1$. Start with $c=0$. If $s^ic=a_i$ for $i<j$, replace $c$ by
\[
 c+d^{j+1}(a_j-s^jc).
\]
The identities $s^jd^{j+1}=\id$ and $s^id^{j+1}=d^js^i$ for $i<j$ show that this corrects the $j$th equation and preserves the previous ones: the latter assertion uses $s^i a_j=s^{j-1}a_i$. This construction is linear and simplicial. Thus each matching map is a split surjection of simplicial abelian groups, hence a Kan fibration. In cosimplicial degree zero the assertion is just the Kan property already noted.

Totalization is right Quillen for the Reedy structure; see the statement immediately before \cite[Proposition 8.1, p.~1034]{FGHBousfield2003}, or check the horn lifting adjunction against the Reedy cofibrations $\Delta^\bullet\otimes\Lambda^k[m]\to\Delta^\bullet\otimes\Delta[m]$. Therefore
\[
 Z_n=\operatorname{Tot}(R^{\bullet+1}K)
\]
is Kan and is a model of $\widehat S^n_2$. In simplicial degree $m$, its set of simplices is the closed compatibility subspace of
\[
 \prod_{q\geq0}
 \operatorname{Hom}_{\mathrm{sSet}}
     (\Delta^q\times\Delta^m,R^{q+1}K).
\]
Each factor is finite: the product of the two standard simplices has finitely many nondegenerate simplices. Give each factor the discrete topology. The displayed product is compact metrizable, its compatibility equations are closed, and all face and degeneracy operators are continuous. Thus $Z_n$ is a simplicial compact metrizable Hausdorff space with underlying Kan complex. This argument supplies an actual strict simplicial model, not just compact topologies on homotopy classes.

\emph{A bounded extension lemma.}
Suppose $Z$ is any pointed Kan complex whose levels have compact metrizable Hausdorff topologies and continuous simplicial operators. For a pointed simplicial set $A$, the set
\[
 H_Z(A)=\operatorname{Hom}_{\mathrm{sSet}_*}(A,Z)
\]
is a closed subspace of the product of the target levels indexed by simplices of $A$, hence is compact Hausdorff. For countable $A$ it is metrizable. Restrictions are continuous.

Fix a pointed simplicial set $X$ and a countable pointed simplicial subset $A\subseteq X$. Let $C_A$ be the image of $H_Z(X)\to H_Z(A)$. This is compact and closed. For $A\subseteq L\subseteq X$ obtained by adjoining finitely many nondegenerate simplices and all their faces, let $C_{L,A}$ be the image of $H_Z(L)\to H_Z(A)$. Then
\begin{equation}\label{eq:G-compact-image-intersection}
 C_A=\bigcap_L C_{L,A}.
\end{equation}
Indeed a compatible assignment on $A$ belongs to every image on the right precisely when every finite set of additional simplex equations can be solved. In the compact product of all the target coordinates indexed by $X$, those equations and the prescribed coordinates on $A$ are closed conditions with the finite intersection property. Compactness gives a simultaneous solution, which is the required map $X\to Z$.

Choose a countable basis of $H_Z(A)$. The basic open sets whose closures miss $C_A$ cover its complement: compact metrizable spaces are regular, so the basis may be chosen to have this property. For each such open $U$, equation \eqref{eq:G-compact-image-intersection} and compactness of $\overline U$ show that finitely many of the sets $C_{L,A}$ already have intersection disjoint from $\overline U$. Their finite union of added simplices gives one countable enlargement $L_U$ with this property. The union of these $L_U$ is a countable enlargement $A'\subseteq X$ such that
\begin{equation}\label{eq:G-one-step-extension}
 \operatorname{im}(H_Z(A')\to H_Z(A))=C_A.
\end{equation}
The same construction can be performed simultaneously for a countable family of such targets by taking the countable union of the enlargements.

Apply this to the countable family
\[
 Z_{n,m}=\operatorname{Map}_*(\Delta^m/\partial\Delta^m,Z_n),
       \qquad n\geq2,\quad m\geq0,
\]
where $m=0$ means $Z_{n,0}=Z_n$. These are Kan models of $\Omega^m Z_n$ and again have compact metrizable levels with continuous operators: each level is a closed subset of a finite product of target levels, because the simplicial sphere smashed with a standard simplex is finite. Starting at any countable $A_0\subseteq X$, iterate \eqref{eq:G-one-step-extension} to obtain countable
\[
 A_0\subseteq A_1\subseteq A_2\subseteq\cdots,
 \qquad A_\infty=\bigcup_i A_i.
\]
Every map $h:A_\infty\to Z_{n,m}$ extends to $X$. To see this, its restriction to $A_i$ extends globally because it is the restriction of its map on $A_{i+1}$. The sets of global extensions of $h|_{A_i}$ are nonempty nested closed subsets of the compact space $H_{Z_{n,m}}(X)$. Their intersection supplies an extension of $h$ itself.

Now assume $X\in\widehat{\mathcal N}_2$. Each such global extension is pointed nullhomotopic, since $\Map_*(X,Z_{n,m})$ is weakly contractible. Its restriction makes $h$ nullhomotopic. Thus
\[
 [A_\infty,\Omega^m Z_n]_* = *
       \quad\text{for all }n\geq2,\ m\geq0.
\]
The $m=0$ assertion makes each $\Map_*(A_\infty,Z_n)$ connected; the other assertions make every homotopy group at its basepoint zero. These mapping spaces are Kan, so they are weakly contractible. Hence $A_\infty\in\widehat{\mathcal N}_2$.

We have proved that every countable simplicial subset of a member of $\widehat{\mathcal N}_2$ is contained in a countable member of that class. Realization converts this to the CW version of B.1: the realization of a countable simplicial set has countably many cells, and every countable collection of cells in a realization is contained in the realization of a countable simplicial subset. Starting with a simplicial model of an arbitrary CW homotopy type and then using the model-invariance argument of B.1 gives
$\widehat{\mathcal N}_2=\widehat{\mathcal N}_2^{\aleph_1}$. This step does not assume that the singular complex of a countable CW complex is countable.

\emph{The rational condition and conclusion.}
The class of connected rationally acyclic spaces has the same countable enlargement property. Given a countable subcomplex of a rationally acyclic $X$, there are only countably many finite rational chains and cycles on it. Adjoin the finite supports of bounding chains in $X$ for all its reduced cycles, and repeat countably. Every cycle in the resulting union occurs at some stage and bounds at the next; the union is therefore rationally acyclic. This includes reduced degree zero and so also ensures connectedness.

If $X\in\Nul$, G.6 says that $X$ is both completed-sphere-null and rationally acyclic. Alternate the two countable enlargement constructions just proved. Their common sequential union belongs to both classes, because both are closed under pointed homotopy colimits. It therefore belongs to $\Nul$, again by G.6. Thus $\Nul=\Nul^{\aleph_1}$. B.2 applies to both null classes, proving the stated accessibility and the explicit generating sets. Cisinski's localization construction used in G.2 now gives the reflector and its pointed derivator in $\mathcal U$, so the suspension and coefficient conclusions of G.3 are unconditional ZFC statements for the full $W_S$.
\end{proof}
